\SourceLeaf{132}
Soll dagegen der Logarithmus der Differenz bestimmt werden, so findet man $\log a - \log b$ in der zweiten oder dritten Spalte, je nachdem, ob diese Differenz kleiner oder größer als $\log 2$ ist. Im ersten Fall, wenn man $\log a - \log b = Q$ setzt, ergibt sich $\log(a-b) = \log b - P = \log a - R$. Im zweiten Fall dagegen, wenn man $\log a - \log b = \log R$ setzt, ist der gesuchte Logarithmus
\[
= \log b + P = \log a - Q.
\]

Es ist sonderbar, dass \textsc{Leonelli} nicht gleich bemerkt hat, wie man den Logarithmus der Differenz auf diese Weise bestimmen kann, wenn die Differenz der Logarithmen kleiner als $\log 2$ ist. In der Schrift selbst gibt er für diesen Fall ein anderes, verwickelteres Verfahren an. Erst der Hinweis auf diesen Mangel im Bericht von \textsc{Delambre} hat ihn auf diese doch so naheliegende Art der Tabellennutzung gebracht. In seinen Bemerkungen zu diesem Bericht erklärt er sie sehr ausführlich.

Wir können \textsc{Leonellis} Gedanken, die logarithmischen Rechnungen durch eine solche Tafel zu erleichtern, insgesamt unsere Zustimmung nicht versagen und halten die tatsächliche Erstellung einer solchen Tafel für wünschenswert. Allem Übrigen, was \textsc{Leonelli} über diesen Gegenstand sagt, können wir jedoch nur wenig Wert beimessen. Seine Erläuterung des Gebrauchs ist für einen so elementaren Gegenstand unnötig weitschweifig. Etwas tiefer gehende Untersuchungen über das Gesetz, nach dem sich die Tafeln und die Differenzen entwickeln, findet man gar nicht, obwohl sie zur Bestimmung des nötigen Umfangs der Tafeln sehr wichtig wären. Es finden sich jedoch einige beiläufige, zum Teil recht verworren formulierte Äußerungen, aus denen man schließen kann, dass \textsc{Leonelli} solche Untersuchungen gar nicht oder ganz fehlerhaft angestellt hat. Dazu gehört z.\,B. die grundfalsche Behauptung, dass, wenn $P$ stets um die Differenz 0,0005 zunimmt und sich dem Wert 6 nähert, die ersten Differenzen von $Q$ nur noch eine oder zwei signifikante Ziffern haben sollen, wenn $Q$ mit 14 Dezimalstellen ausgedrückt wird. Eine einfache Rechnung zeigt, dass diese Differenzen bis dahin mindestens fünf signifikante Ziffern behalten. Noch weniger verschwinden sie, wenn $P$ den Wert 6 übersteigt, sondern erst für $P = 10,7$ werden sie bis auf eine Einheit in der vierzehnten Dezimalstelle abge\SourceLeaf{133}%
